% ==============================================================================
% FFGFT: Document 244 (Chapter content)
% Title: Hilbert Space Comparison: FFGFT L²(T⁴) and UMF \ell ²(P_N)
% Author: Johann Pascher
% Date: May 2026
% Internal analysis — basis for response to Marco Gericke (UMF)
% ==============================================================================

\section*{Abstract}

Marco Gericke (Universal Model Framework, UMF) has proposed a formal architecture
in which quantum mechanics, special relativity, and fermion mass emerge from a
prime-indexed Hilbert space $\ell^2(P_N)$ and its reciprocal interior
$\ell^2(P_N^{-1})$, coupled by a bivector $B_{OI}$. This document analyses the
structural relationship between UMF and FFGFT. The central result is:
FFGFT's $L^2(T^4)$ is the more complete Hilbert space — it contains UMF's
$\ell^2(P_N)$ as a proper subspace. However, UMF's prime-indexed structure offers
genuine computational advantages in specific domains, particularly prime
factorisation and number-theoretic problems. The two frameworks are therefore not in competition but stand in a relation of \textbf{mathematical containment (subset relation) and physical complementarity}: UMF's Hilbert space is a proper subspace of FFGFT's, but UMF's prime-indexed structure, bivector $B_{OI}$, and activation operator $K_\text{Kriger}$ provide formal tools that FFGFT currently lacks.

\subsection*{Notation}

\begin{itemize}
  \item $P_N = \{p_1, p_2, \dots, p_N\}$ with $p_1=2, p_2=3, p_3=5, \dots$ (first $N$ primes)
  \item $T^4 = (\mathbb{R}/2\pi\mathbb{Z})^4$ (4-torus, FFGFT state space)
  \item $B_{OI} = \Pi_\text{out} \wedge \Pi_\text{in} \in \Lambda^2(H_\text{out} \oplus H_\text{in})$
    (UMF bivector, see \S\ref{sec:bivector})
  \item $\xi = 4/30000$ (FFGFT dimensionless winding parameter)
  \item $K_\text{Kriger}$: UMF activation operator (see \S\ref{sec:kriger})
  \item $\xi_\text{UIFT}$: dimensionless parameter of the UIFT framework (Doc.~013), related to $\xi_\text{FFGFT}$ via the bridge factor $414{,}684$
\end{itemize}

% ==============================================================================
\section*{List of Symbols}
% ==============================================================================

\begin{center}
\renewcommand{\arraystretch}{1.4}
{\small
\begin{tabular}{p{3.5cm}p{9cm}}
\toprule
\textbf{Symbol} & \textbf{Meaning} \\
\midrule
$\xi = 4/30000$ & FFGFT dimensionless winding parameter (geometrically fixed, $\kappa=7$) \\
$\xi_\text{FFGFT}$ & Same as $\xi$; used when distinguishing from $\xi_\text{UIFT}$ \\
$\xi_\text{UIFT}$ & Dimensionless parameter of UIFT: $k_BT_\text{CMB}\ln 2/(m_ec^2)$ (Doc.~013) \\
$L^2(T^4)$ & Hilbert space of square-integrable functions on the 4-torus $T^4$ \\
$T^4$ & 4-dimensional torus $(\mathbb{R}/2\pi\mathbb{Z})^4$; FFGFT state space \\
$\ell^2(P_N)$ & UMF Hilbert space: square-summable sequences over first $N$ primes \\
$P_N$ & Set of first $N$ prime numbers $\{2,3,5,\ldots,p_N\}$ \\
$w_p = \ln p/\sqrt{p}$ & Prime weight in UMF inner product \\
$B_{OI}$ & UMF bivector: $\Pi_\text{out} \wedge \Pi_\text{in}$; couples the two spaces \\
$\Pi_\text{out}$ & UMF Prime Boundary $= \ell^2(P_N)$ (exterior, metric domain) \\
$\Pi_\text{in}$ & UMF Reciprocal Interior $= \ell^2(P_N^{-1}) \oplus \ell^2(Q_{ZN})$ (mass domain) \\
$K_\text{Kriger}$ & UMF activation operator; determines which interface information generates curvature \\
$\chi(x)$ & Dimensionless activation parameter: $k_BT(x)\ln2\cdot\tau_t(x)/(\hbar/4)$ \\
$\chi = 1$ & Landauer threshold: boundary between inactive and active structure \\
$D_P$ & UMF chiral Dirac operator on $\ell^2(P_N)\otimes\mathbb{C}^2$ \\
$\tau\cdot\mu$ & Coherence-time -- information-mass product; $= \hbar/c^2$ (both frameworks) \\
$\varepsilon \approx 0.001$ & FFGFT non-closure of $T^4$; generative structural residual \\
$\theta_4$ & Continuous fourth-torus phase coordinate; absent from $\ell^2(P_N)$ \\
$n_\theta, n_\phi$ & FFGFT winding numbers in two independent torus directions \\
$\iota$ & Embedding $\ell^2(P_N) \hookrightarrow L^2(T^4)$ (containment map) \\
\bottomrule
\end{tabular}
}
\end{center}


\clearpage

% ==============================================================================
\section{The Two Hilbert Spaces}
% ==============================================================================

\subsection{FFGFT: \texorpdfstring{$L^2(T^4)$}{L2(T4)}}

FFGFT's state space is $L^2(T^4)$ — the space of square-integrable functions
on the 4D torus $T^4 = (\mathbb{R}/2\pi\mathbb{Z})^4$. The natural orthonormal
basis consists of Fourier modes:
\begin{equation}
  e_{\mathbf{n}}(\symbf{\theta}) = \frac{1}{(2\pi)^2}
  e^{i(n_1\theta_1 + n_2\theta_2 + n_3\theta_3 + n_4\theta_4)},
  \qquad \mathbf{n} \in \mathbb{Z}^4
  \tag{244.1}
\end{equation}
This space is:
\begin{itemize}
  \item \textbf{Complete}: every Cauchy sequence converges in $L^2(T^4)$
  \item \textbf{Separable}: the Fourier basis $\{e_{\mathbf{n}}\}_{\mathbf{n}\in\mathbb{Z}^4}$
    is countable and dense
  \item \textbf{Continuous}: admits continuous winding numbers and phase structure
  \item \textbf{Indexed by} $\mathbb{Z}^4$: all integer winding number combinations,
    including prime-indexed ones as a special case
\end{itemize}

The single parameter $\xi = 4/30000$ determines which winding modes are
dynamically preferred — but all modes are structurally present in $L^2(T^4)$.

\subsection{UMF: \texorpdfstring{$\ell^2(P_N)$}{ell2(P\_N)}}

UMF's primary space is $\ell^2(P_N)$ — the space of square-summable sequences
indexed by the first $N$ prime numbers, with inner product:
\begin{equation}
  \langle f | g \rangle_{P_N} = \sum_{p \leq p_N} w_p\, \overline{f(p)}\, g(p),
  \qquad w_p = \frac{\ln p}{\sqrt{p}}
  \tag{244.2}
\end{equation}
This space is:
\begin{itemize}
  \item \textbf{Complete} for fixed finite $N$: $\ell^2(P_N) \cong \mathbb{C}^N$
  \item \textbf{Discrete}: no continuous phase structure
  \item \textbf{Indexed by primes only}: $\{2, 3, 5, 7, 11, \ldots, p_N\}$
  \item \textbf{Weight-divergent} as $N\to\infty$: $\sum_p w_p = \sum_p \ln p/\sqrt{p}$
    diverges (by comparison with $\sum 1/\sqrt{p}$)
\end{itemize}

% ==============================================================================
\section{Containment Relation}
% ==============================================================================

\subsection{\texorpdfstring{$\ell^2(P_N)$}{ell2(P\_N)} is a Subspace of \texorpdfstring{$L^2(T^4)$}{L2(T4)}}

The prime-indexed basis vectors of UMF correspond to specific Fourier modes
of $L^2(T^4)$. Consider the embedding:
\begin{equation}
  \iota: \ell^2(P_N) \hookrightarrow L^2(T^4),
  \qquad f(p) \mapsto f(p)\cdot e_{(p,0,0,0)}(\symbf{\theta})
  \tag{244.3}
\end{equation}
This maps each prime-indexed coefficient to a Fourier mode with winding number
$(p, 0, 0, 0)$ in the first torus direction. The embedding is isometric up to
normalisation and injective. Therefore:
\begin{equation}
  \boxed{\ell^2(P_N) \hookrightarrow L^2(T^4), \quad \text{proper subspace}}
  \tag{244.4}
\end{equation}

$L^2(T^4)$ contains far more: all integer winding numbers in all four directions,
all non-prime modes, all mixed-direction modes, and the continuous phase structure
$\theta_4$ which is absent from $\ell^2(P_N)$ entirely.

\subsection{What UMF's Space Cannot Represent}

The following FFGFT structures have no counterpart in $\ell^2(P_N)$:

\begin{center}
\renewcommand{\arraystretch}{1.4}
{\small
\begin{tabular}{p{5cm}p{4cm}p{4cm}}
\toprule
\textbf{FFGFT Structure} & \textbf{In $L^2(T^4)$} & \textbf{In $\ell^2(P_N)$} \\
\midrule
Continuous phase $\theta_4$ & $\checkmark$ (full range) & $\times$ (absent) \\
Non-prime winding modes & $\checkmark$ & $\times$ \\
Mixed-direction winding $(n_1,n_2,n_3,n_4)$ & $\checkmark$ & $\times$ \\
Non-closure $\varepsilon$ (winding residual) & $\checkmark$ & $\times$ \\
Z3$^3$ topological constraint & $\checkmark$ & partially \\
4$\pi$ spinor periodicity & $\checkmark$ & $\times$ \\
Lorentz structure (from $T^4$ geometry) & $\checkmark$ & derived via RG \\
\bottomrule
\end{tabular}
}
\end{center}

\subsection{The Missing Object: UMF Bivector \texorpdfstring{$B_{OI}$}{B\_OI}}
\label{sec:bivector}

Marco Gericke's central formal contribution is not $\ell^2(P_N)$ alone but the
wedge product of the two spaces:
\begin{equation}
  B_{OI} = \Pi_\text{out} \wedge \Pi_\text{in}
  \;\in\; \Lambda^2(H_\text{out} \oplus H_\text{in})
  \tag{244.4b}
\end{equation}
where $\Pi_\text{out} = \ell^2(P_N)$ and
$\Pi_\text{in} = \ell^2(P_N^{-1}) \oplus \ell^2(Q_{ZN})$.

$B_{OI}$ has three crucial properties:
\begin{enumerate}
  \item \textbf{Projection-stable}: neither $\pi_1$ (metric projection) nor
    $\pi_2$ (spinor/mass projection) alone contains it, but both projections
    preserve its existence as a structural invariant.
  \item \textbf{Generates observable geometry}:
    $g_{\mu\nu}^\text{phys} = g_{\mu\nu}^{(0)} + \lambda_B\,
    B_{\mu\alpha}^\text{act} B_{\nu}^{\text{act}\,\alpha}$
  \item \textbf{Not derivable from either space alone}: it is the coupling
    between $\Pi_\text{out}$ and $\Pi_\text{in}$, not a state within either.
\end{enumerate}

This is the formal realisation of what Stefaan Vossen called the GIA
projection-stable invariant, and what FFGFT described verbally as
``mutual constitution''. $B_{OI}$ is \emph{not} an element of $L^2(T^4)$
— it is a relation \emph{between} spaces. The containment
$\ell^2(P_N) \hookrightarrow L^2(T^4)$ holds for the spaces themselves;
$B_{OI}$ is a genuinely new object that FFGFT currently has no algebraic
counterpart for.

% ==============================================================================
\section{Where UMF Offers Genuine Added Value}
% ==============================================================================

The containment relation does not make UMF redundant. Three domains where
the prime-indexed structure offers advantages:

\subsection{Prime Factorisation and Number Theory (Doc 075, 076)}

FFGFT Doc 075 treats RSA factorisation via T0-field dynamics — period-finding
as wave propagation rather than quantum superposition. The natural computational
basis for this problem is prime-indexed. UMF's $\ell^2(P_N)$ with the
prime-weighted inner product $w_p = \ln p/\sqrt{p}$ is the \emph{natural}
Hilbert space for number-theoretic computations: the weights encode the
prime number theorem density $\pi(x) \sim x/\ln x$ directly.

In $L^2(T^4)$, prime modes are present but not distinguished. In $\ell^2(P_N)$,
they are the entire basis. For factorisation problems, UMF's representation
is computationally more efficient.

\subsection{The Dirac Operator \texorpdfstring{$D_P$}{D\_P} as Computable Spectrum}

UMF's chiral Dirac operator on $\ell^2(P_N) \otimes \mathbb{C}^2$:
\begin{equation}
  D_P = \begin{pmatrix} 0 & A \\ A^\dagger & 0 \end{pmatrix},
  \quad A_{pq} = w_p\,\delta_{pq}
  \tag{244.5}
\end{equation}
has a discrete, explicitly computable spectrum. The spectrum of the
corresponding operator on $L^2(T^4)$ requires solving the full torus
eigenvalue problem — substantially harder. For mass-spectrum calculations
in the lepton sector, both approaches give $\sim 1\%$ agreement with PDG,
but UMF's computation is more direct.

\subsection{Independent Confirmation of Lorentz Invariance}

UMF derives Lorentz invariance as the IR fixed point of the anisotropy RG
flow $u(\mu) \to u^* = 1$. FFGFT derives it from the $T^4$ winding geometry
and the constraint $\tilde{T}\cdot m = 1$. These are structurally independent
derivations arriving at the same result. The independence strengthens both:
if two different mathematical starting points require Lorentz invariance,
the result is more robust.

\subsection{The Activation Operator \texorpdfstring{$K_\text{Kriger}$}{K\_Kriger}}
\label{sec:kriger}

UMF defines the dimensionless activation parameter:
\begin{equation}
  \chi(x) = \frac{k_B T(x)\ln 2 \cdot \tau_t(x)}{\hbar/4}
  \tag{244.6a}
\end{equation}
and the soft activation gate:
\begin{equation}
  A(x) = \frac{1}{1+e^{-s(\chi(x)-1)}}
  \tag{244.6b}
\end{equation}
Physical geometry receives contributions only from the activated fraction:
\begin{equation}
  \Phi_\text{grav} = A(x) \cdot \Phi_\text{interface}
  \tag{244.6c}
\end{equation}

\textbf{Key insight:} The threshold $\chi = 1$ is exactly the Landauer bound.
For $\chi < 1$, the bivector $B_{OI}$ exists mathematically but generates no
curvature — it is structurally present but physically inactive. This formalises
what FFGFT has described qualitatively as the boundary between geometrically
present and observationally active structure. FFGFT has no algebraic analogue
of $K_\text{Kriger}$. It is UMF's genuine formal contribution, independent
of the containment relation.

% ==============================================================================
\section{\texorpdfstring{The $\tau\cdot\mu$ Identity}{The tau.mu Identity}: Independent Derivation}
% ==============================================================================

Both frameworks arrive at the same identity independently:
\begin{equation}
  \tau \cdot \mu = \frac{\hbar}{c^2}
  \tag{244.7}
\end{equation}
where $\tau = \hbar/(k_B T_\text{CMB}\ln 2)$ and $\mu = k_B T_\text{CMB}\ln 2/c^2$.

In FFGFT this is the coherence time -- information-mass product at the CMB scale.
In UMF this is the statement that the Kriger activation threshold, evaluated at
$T_\text{CMB}$, equals the Compton-scale information mass. Two different
derivational paths, identical result. The agreement is not coincidental:
it reflects the bridge formula $\xi_\text{FFGFT}/\xi_\text{UIFT} \approx 414{,}684$
(Doc 190 R13), which connects the same CMB temperature to both frameworks
through different normalisations.

% ==============================================================================
\section{The Non-Redundancy Principle}
% ==============================================================================

The most useful framework is not the most complete one but the one that
derives the most from the fewest assumptions. By this criterion:

\textbf{Table 1 — Technical comparison:}

\begin{center}
\renewcommand{\arraystretch}{1.4}
{\small
\begin{tabular}{p{4cm}p{4cm}p{4cm}}
\toprule
\textbf{Criterion} & \textbf{FFGFT} & \textbf{UMF} \\
\midrule
Free parameters & 1 ($\xi$) & 5 (Distinction Logic axioms + $N$) \\
Hilbert space & $L^2(T^4)$ (complete) & $\ell^2(P_N) \subset L^2(T^4)$ \\
Lorentz derivation & from $T^4$ geometry & from RG flow (independent) \\
Mass spectrum & from winding recursion & from $D_P$ spectrum (compact) \\
Activation boundary & qualitative & $K_\text{Kriger}$ (explicit) \\
Factorisation & wave propagation & natural prime basis \\
\bottomrule
\end{tabular}
}
\end{center}

\textbf{Table 2 — Ontological / structural comparison:}

\begin{center}
\renewcommand{\arraystretch}{1.4}
{\small
\begin{tabular}{p{4cm}p{4cm}p{4cm}}
\toprule
\textbf{Criterion} & \textbf{FFGFT} & \textbf{UMF} \\
\midrule
Primitive & Geometry ($T^4$) & Distinction Logic \\
Bivector $B_{OI}$ & not present & $\Pi_\text{out} \wedge \Pi_\text{in}$ \\
``Mutual constitution'' & described verbally & algebraically realised \\
Ontological commitment & geometric primacy & distinction-logic primacy \\
\bottomrule
\end{tabular}
}
\end{center}

The two frameworks occupy different positions in the projection space $\mathcal{M}$
(Stefaan Vossen's comparative architecture): they share projection-stable
invariants (Lorentz structure, $\tau\cdot\mu$, mass spectrum) while differing
in generative ontology and computational representation.

% ==============================================================================
\section{Conclusion: Complementarity without Redundancy}
% ==============================================================================

The analysis yields five results that are independent of any particular
exchange or context and follow directly from the structural comparison above.

\begin{enumerate}

  \item \textbf{Containment.}
    $\ell^2(P_N) \hookrightarrow L^2(T^4)$: UMF's Hilbert space is a proper
    subspace of FFGFT's. $L^2(T^4)$ is the structurally more complete space.
    This is a mathematical fact, not a value judgement about either framework.

  \item \textbf{Non-redundancy.}
    The containment relation does not make UMF redundant. UMF offers
    computational advantages in prime-indexed domains, an explicit activation
    operator ($K_\text{Kriger}$), an independent derivation of Lorentz
    invariance via RG flow, and a computable Dirac spectrum. These are genuine
    contributions not present in $L^2(T^4)$.

  \item \textbf{Structural alignment.}
    The \texorpdfstring{$\tau\cdot\mu$}{tau.mu} identity $= \hbar/c^2$
    is derived independently in both frameworks and confirms structural
    convergence at the CMB scale (Doc.~013).

  \item \textbf{The bivector as a new object.}
    UMF's $B_{OI} = \Pi_\text{out} \wedge \Pi_\text{in}$ and the activation
    operator $K_\text{Kriger}$ are formal objects FFGFT currently lacks.
    They are relations \emph{between} Hilbert spaces, not states within one
    space, and are therefore not contained in $L^2(T^4)$. A natural next step
    is to map $\xi_\text{FFGFT}$ to the UMF compactification radius $R_\chi$
    via the bridge factor $\xi_\text{FFGFT}/\xi_\text{UIFT} \approx 414{,}684$
    (Doc.~013), and to complete the UMF entry in the comparative projection
    schema (Vossen 2026).

  \item \textbf{Ontological neutrality.}
    No ontological priority claim is required by this analysis. UMF derives
    $\ell^2(P_N)$ from distinction logic; FFGFT derives $L^2(T^4)$ from
    geometry. The subspace relation holds regardless of which generative
    account one adopts.

\end{enumerate}
